\documentclass[12pt]{article}
\usepackage{it_hw}
\usepackage{amssymb} % \blacksquare and other AMS symbols

\newcommand{\E}{\field{E}}
\input{macro}
% macro.tex redefines \E as a one-argument operator \E{x}=E[x]; we use \E as the
% bare symbol with explicit subscripts/brackets, so restore the no-argument form.
\renewcommand{\E}{\field{E}}
% Some minimal TeX installs lack the slanted-Helvetica footer metric (phvro7t);
% fall back to the default family for the page footer if it is unavailable.
\IfFileExists{phvro7t.tfm}{}{\renewcommand{\hffont}{\fontsize{9}{11}\selectfont}}
% Mode toggle: compile with  \def\ExamMode{}  for the blank exam (final26.pdf);
% the default (no ExamMode) produces the solutions (final26sol.pdf).
\ifdefined\ExamMode
  \includeversion{problem}
  \excludeversion{solution}
\else
  \excludeversion{problem}
  \includeversion{solution}
\fi

\begin{document}

\begin{center}
{\bf Information Theory and Its Application in Deep Learning: Final \begin{solution}Solutions\end{solution}} \\ \medskip
\begin{problem}June 2026 \end{problem}
\end{center}

\begin{problem}
\noindent
\begin{itemize}
\item You can use previous parts of a problem even if you did not solve them.
\item Logarithms are natural (nats) unless stated otherwise; for the rate--distortion problem use base $2$ ($\log = \log_2$, rates in bits).
\item $D(\cdot\|\cdot)$ denotes the KL divergence and $h(\cdot)$ the differential entropy.
\end{itemize}
\LeftFoot{Final}
\end{problem}
\begin{solution}
\LeftFoot{Final Solution}
\end{solution}

\begin{problem}

    \vskip .3in
    {\bf ID : \underline{\hspace{6cm}}}\\
    \vskip .3in
    {\bf Name : \underline{\hspace{5.5cm}}}\\
    ~
    \vskip 2in
    \begin{table}[h]
        \Large
        \begin{tabular}{c|c}
           1 &\textcolor{white}{aaaaaaaa}/~30\\
           2 &\textcolor{white}{aaaaaaaa}/~35 \\
           3 &\textcolor{white}{aaaaaaaa}/~35 \\
            \hline
           Total &  \textcolor{white}{aaaaaaaa}/~100 \\
        \end{tabular}
    \end{table}
    \newpage
    {\bf extra sheet}
    \newpage
\end{problem}

\begin{enumerate}

% Problem 1: Differential Entropy and Maximum Entropy
%%%%%%%%%%%%%%%%%%%%%%%
\item {\bf Differential Entropy and Maximum Entropy} {(\it 30 points)}\\
Let $X$ be a continuous random variable. Recall $h(X) = -\int f(x)\log f(x)\,dx$ (nats).
\begin{enumerate}
\item Let $X \sim \mathrm{Exp}(\lambda)$, i.e.~$f(x) = \lambda e^{-\lambda x}$ for $x \ge 0$. Compute $h(X)$. {(\it 10 points)}
\item Let $Y = cX$ for a constant $c > 0$. Find $h(Y)$ in two ways: (i) from the scaling rule, and (ii) by identifying the distribution of $Y$. {(\it 10 points)}
\item Show that among \emph{all} densities supported on $(0,\infty)$ with mean $\E[X] = \mu$, the exponential density $f(x) = \frac1\mu e^{-x/\mu}$ maximizes $h$. {(\it 10 points)}
\end{enumerate}

%Solution
%%%%%%%%%%%%%%%%%%%%%%%
\vskip .2in
\begin{solution}{\bf Solution: Differential Entropy and Maximum Entropy.}
\begin{enumerate}
\item Take the logarithm of the density first: $\log f(x) = \log\lambda - \lambda x$. Substituting into the definition and splitting the integral,
\[
\begin{aligned}
    h(X) &= -\int_0^\infty \lambda e^{-\lambda x}\big(\log\lambda - \lambda x\big)\,dx\\
         &= -\log\lambda\underbrace{\int_0^\infty \lambda e^{-\lambda x}\,dx}_{=\,1}
           \;+\; \lambda\underbrace{\int_0^\infty x\,\lambda e^{-\lambda x}\,dx}_{=\,\E[X]}.
\end{aligned}
\]
The first integral is $1$ because $f$ is a density; the second is $\E[X]=\tfrac1\lambda$. Hence
\[
    h(X) = -\log\lambda + \lambda\cdot\tfrac1\lambda = 1 - \log\lambda.
\]
\item {\bf (i) Scaling rule.} For $Y=cX$ with $c>0$ the density transforms as $f_Y(y)=\tfrac1c f_X(y/c)$, so
\[
    h(Y) = -\int f_Y(y)\log f_Y(y)\,dy
         = -\int \tfrac1c f_X(y/c)\big[\log f_X(y/c) - \log c\big]\,dy.
\]
Substituting $x=y/c$ (so $dy = c\,dx$) gives $h(Y) = -\int f_X(x)\big[\log f_X(x)-\log c\big]\,dx = h(X)+\log c$. Therefore $h(Y) = 1 - \log\lambda + \log c$.\\
{\bf (ii) Identify the law.} If $X\sim\mathrm{Exp}(\lambda)$ then $Y=cX$ has $\Pr(Y\le y)=\Pr(X\le y/c)=1-e^{-\lambda y/c}$, i.e.~$Y\sim\mathrm{Exp}(\lambda/c)$. By part (a) with rate $\lambda/c$, $h(Y) = 1 - \log(\lambda/c) = 1 - \log\lambda + \log c$. The two answers agree.
\item Let $g$ be any density on $(0,\infty)$ with $\int_0^\infty x\,g(x)\,dx = \mu$, and let $f(x)=\tfrac1\mu e^{-x/\mu}$ be the exponential with the same mean. Start from non-negativity of KL and expand the integrand:
\[
    0 \le D(g\|f) = \int g\log\frac{g}{f}
      = \int g\log g - \int g\log f
      = -h(g) - \int g(x)\log f(x)\,dx.
\]
The key point is that $\log f$ is \emph{linear} in $x$: $\log f(x) = -\log\mu - x/\mu$, so its $g$-average depends on $g$ only through the mean,
\[
    \int g(x)\log f(x)\,dx = -\log\mu - \frac1\mu\int_0^\infty x\,g(x)\,dx = -\log\mu - 1,
\]
which is the same value it takes under $f$ itself. Substituting back, $0 \le -h(g) + \log\mu + 1$, i.e.
\[
    h(g) \le 1 + \log\mu .
\]
The bound is attained by $f$, since $h(f) = 1 - \log(1/\mu) = 1 + \log\mu$ (part (a) with $\lambda=1/\mu$). Hence the exponential maximizes $h$ over this class, with maximum entropy $1+\log\mu$, and equality holds iff $D(g\|f)=0$, i.e.~$g=f$.
\end{enumerate}
\end{solution}

\begin{problem}
    \newpage~\newpage
\end{problem}

% Problem 2: Rate-Distortion of an Independent Gaussian Pair
%%%%%%%%%%%%%%%%%%%%%%%
\vskip .3in
\item {\bf Rate--Distortion of an Independent Gaussian Pair} {(\it 35 points)}\\
Consider a vector source $\bX = (X_1, X_2)$ with \emph{independent} components $X_i \sim \mathcal{N}(0,\sigma_i^2)$, and squared-error distortion $d(\mathbf{x},\hat{\mathbf{x}}) = \sum_{i=1}^2 (x_i - \hat x_i)^2$. Work with the information rate--distortion function
\[
    R(D) = \min_{p(\hat{\mathbf x}\mid \mathbf x):\,\E\|\bX-\bXh\|^2\le D} I(\bX;\bXh),
\]
and write $R_i$ for the single-component function $R_i(D_i) = \frac12\log\frac{\sigma_i^2}{D_i}$ ($0<D_i\le\sigma_i^2$). Rates are in bits, $\log=\log_2$. The goal is to prove
\[
    R(D) = \min_{\substack{D_1,D_2\ge 0\\ D_1+D_2\le D}} \big(R_1(D_1)+R_2(D_2)\big).
\]
\begin{enumerate}
\item {\bf Decoupling lemma.} For \emph{any} joint distribution of $(\bX,\bXh)$, show
\[
    I(\bX;\bXh) \;\ge\; I(X_1;\hat X_1) + I(X_2;\hat X_2).
\]
\emph{Hint: write $I = h(\bX) - h(\bX\mid\bXh)$; use independence on the first term, the chain rule on the second, and ``conditioning reduces entropy.''} {(\it 10 points)}
\item {\bf Converse.} Show that
\[
    R(D) \;\ge\; \min_{D_1+D_2\le D}\big(R_1(D_1)+R_2(D_2)\big).
\]
\emph{Hint: take any test channel feasible for $R(D)$; set $D_i=\E[(X_i-\hat X_i)^2]$ so $D_1+D_2\le D$. Apply part (a), then bound each $I(X_i;\hat X_i)$ below using the definition of the single-letter $R_i$.} {(\it 10 points)}
\item {\bf Achievability.} Show the matching upper bound
\[
    R(D) \;\le\; \min_{D_1+D_2\le D}\big(R_1(D_1)+R_2(D_2)\big),
\]
so that equality holds.
\emph{Hint: fix any split $(D_1,D_2)$ with $D_1+D_2\le D$ and build a \emph{product} test channel $p(\hat{\mathbf x}\mid\mathbf x)=\prod_i p_i^\star(\hat x_i\mid x_i)$ from the per-component optimizers. Compute its mutual information and its distortion.} {(\it 15 points)}
\end{enumerate}

%Solution
%%%%%%%%%%%%%%%%%%%%%%%
\vskip .2in
\begin{solution}{\bf Solution: Rate--Distortion of an Independent Gaussian Pair.}
\begin{enumerate}
\item {\bf (Decoupling lemma.)} Write the mutual information in differential-entropy form, $I(\bX;\bXh)=h(\bX)-h(\bX\mid\bXh)$, and treat the two terms separately.

\emph{First term.} Since $X_1\perp X_2$, the joint density factorizes and the entropy is additive: $h(\bX)=h(X_1)+h(X_2)$.

\emph{Second term.} By the chain rule for differential entropy,
\[
    h(\bX\mid\bXh) = h(X_1\mid\bXh) + h(X_2\mid X_1,\bXh).
\]
Conditioning on more variables can only reduce differential entropy. Dropping variables from each conditioning set (keep only $\hat X_1$ in the first term, only $\hat X_2$ in the second) gives the upper bound
\[
    h(X_1\mid\bXh) \le h(X_1\mid\hat X_1),
    \qquad
    h(X_2\mid X_1,\bXh) \le h(X_2\mid\hat X_2),
\]
because $\bXh=(\hat X_1,\hat X_2)$ contains $\hat X_1$, and $(X_1,\bXh)$ contains $\hat X_2$. Adding the two, $h(\bX\mid\bXh)\le h(X_1\mid\hat X_1)+h(X_2\mid\hat X_2)$. Combining the two terms,
\[
    I(\bX;\bXh) = h(\bX)-h(\bX\mid\bXh)
    \ge \sum_{i=1}^2\big(h(X_i)-h(X_i\mid\hat X_i)\big) = \sum_{i=1}^2 I(X_i;\hat X_i). \qquad\blacksquare
\]
\item {\bf (Converse.)} Let $p(\hat{\mathbf x}\mid\mathbf x)$ be \emph{any} test channel feasible for $R(D)$, and define $D_i:=\E[(X_i-\hat X_i)^2]$ for the distortion it induces on each component. Summing, $D_1+D_2=\E\|\bX-\bXh\|^2\le D$, so $(D_1,D_2)$ is a feasible split. The channel's marginal $X_i\to\hat X_i$ is itself a single-letter test channel achieving distortion $D_i$, and the scalar rate--distortion function $R_i(D_i)$ is by definition the \emph{smallest} mutual information over all such channels; therefore $I(X_i;\hat X_i)\ge R_i(D_i)$. Chaining this with part (a),
\[
    I(\bX;\bXh)
    \;\overset{(a)}{\ge}\; \sum_i I(X_i;\hat X_i)
    \;\ge\; \sum_i R_i(D_i)
    \;\ge\; \min_{D_1+D_2\le D}\sum_i R_i(D_i).
\]
The last quantity is a constant independent of the chosen channel. Taking the minimum of the left-hand side over all feasible channels --- which by definition is $R(D)$ --- preserves the inequality:
\[
    R(D) \ge \min_{D_1+D_2\le D}\big(R_1(D_1)+R_2(D_2)\big). \qquad\blacksquare
\]
\item {\bf (Achievability.)} Fix any split $(D_1,D_2)$ with $D_1+D_2\le D$. For each component let $p_i^\star(\hat x_i\mid x_i)$ be the optimal scalar Gaussian test channel attaining $I(X_i;\hat X_i)=R_i(D_i)$ at distortion $D_i$, and form the \emph{product} channel
\[
    p(\hat{\mathbf x}\mid\mathbf x)=p_1^\star(\hat x_1\mid x_1)\,p_2^\star(\hat x_2\mid x_2).
\]
Because both the source ($X_1\perp X_2$) and the channel factorize, the two pairs $(X_1,\hat X_1)$ and $(X_2,\hat X_2)$ are mutually independent, so the joint law satisfies $p(\mathbf x,\hat{\mathbf x})=\prod_i p(x_i,\hat x_i)$. Mutual information of independent pairs is additive,
\[
\begin{aligned}
    I(\bX;\bXh)&=\E\log\frac{p(\bX,\bXh)}{p(\bX)\,p(\bXh)}
    =\sum_i \E\log\frac{p(X_i,\hat X_i)}{p(X_i)\,p(\hat X_i)}\\
    &=\sum_i I(X_i;\hat X_i)=\sum_i R_i(D_i),
\end{aligned}
\]
and the distortion adds by linearity of expectation, $\E\|\bX-\bXh\|^2=\sum_i\E[(X_i-\hat X_i)^2]=D_1+D_2\le D$. This product channel is therefore feasible for $R(D)$, so $R(D)\le \sum_i R_i(D_i)$. Since this holds for \emph{every} feasible split, we may minimize the right-hand side:
\[
    R(D)\le \min_{D_1+D_2\le D}\sum_i R_i(D_i).
\]
Together with the converse in part (b), the two bounds coincide:
\[
    R(D) = \min_{D_1+D_2\le D}\big(R_1(D_1)+R_2(D_2)\big). \qquad\blacksquare
\]
Solving this convex program (Lagrange multipliers) yields the \emph{reverse water-filling} allocation $D_i=\min(\theta,\sigma_i^2)$ --- but that step is not required here.
\end{enumerate}
\end{solution}

\begin{problem}
    \newpage~\newpage
\end{problem}

% Problem 3: The chi-squared Divergence
%%%%%%%%%%%%%%%%%%%%%%%
\vskip .3in
\item {\bf The $\chi^2$-Divergence: From Estimator Variance to the Evidence} {(\it 35 points)}\\
The $\chi^2$-divergence is $\chi^2(P\|Q)=\int\frac{(p-q)^2}{q}=\int\frac{p^2}{q}-1$. After establishing its basic properties, you will use it twice: once to lower-bound the variance of an estimator, and once to upper-bound an intractable log-evidence.
\begin{enumerate}
\item {\bf Foundations.} Show $\chi^2(P\|Q)=\mathrm{Var}_Q\!\big(\frac{dP}{dQ}\big)$, and that $D(P\|Q)\le\chi^2(P\|Q)$. \emph{Hint: $\log t\le t-1$.} {(\it 10 points)}
\item {\bf Hammersley--Chapman--Robbins.} For any statistic $T$ with $\E_P[T]=a$, $\E_Q[T]=b$, prove
\[
    \mathrm{Var}_Q(T) \;\ge\; \frac{(a-b)^2}{\chi^2(P\|Q)}.
\]
Then, for a parametric family $\{P_\theta\}$ with an unbiased estimator of $\theta$, take $Q=P_\theta$, $P=P_{\theta+h}$ and let $h\to0$ to recover the Cram\'er--Rao bound $\mathrm{Var}_\theta(T)\ge 1/I(\theta)$. \emph{Hint: express $a-b$ as a covariance under $Q$ and apply Cauchy--Schwarz.} {(\it 15 points)}
\item {\bf A $\chi^2$ upper bound on the evidence (CUBO).} For a latent-variable model $p(x)=\int p(x,z)\,dz$, a proposal $q(z\mid x)$, and weight $w=\frac{p(x,z)}{q(z\mid x)}$, define $\mathrm{CUBO}_2=\frac12\log\E_{q(z\mid x)}[w^2]$. Show $\log p(x)\le\mathrm{CUBO}_2$, that the gap is exactly $\frac12\log\big(1+\chi^2(p(z\mid x)\|q(z\mid x))\big)$, and hence $\mathrm{ELBO}\le\log p(x)\le\mathrm{CUBO}_2$. \emph{Hint: $p(x)=\E_q[w]$ and $(\E_q[w])^2\le\E_q[w^2]$; then write $w=p(x)\,\frac{p(z\mid x)}{q(z\mid x)}$.} {(\it 10 points)}
\end{enumerate}

%Solution
%%%%%%%%%%%%%%%%%%%%%%%
\vskip .2in
\begin{solution}{\bf Solution: The $\chi^2$-Divergence.}
\begin{enumerate}
\item {\bf (Foundations.)} Let $W=\frac{dP}{dQ}=p/q$ be the likelihood ratio. Its first two $Q$-moments are
\[
    \E_Q[W]=\int q\,\frac pq = \int p = 1,
    \qquad
    \E_Q[W^2]=\int q\,\Big(\frac pq\Big)^2=\int\frac{p^2}{q}.
\]
Hence $\mathrm{Var}_Q(W)=\E_Q[W^2]-(\E_Q[W])^2=\int\frac{p^2}{q}-1=\chi^2(P\|Q)$, the first claim. For the comparison with KL, apply $\log t\le t-1$ with $t=p/q$ pointwise:
\[
    D(P\|Q)=\int p\log\frac pq \;\le\; \int p\Big(\frac pq-1\Big)=\int\frac{p^2}{q}-\int p=\int\frac{p^2}{q}-1=\chi^2(P\|Q).
\]
\item {\bf (Hammersley--Chapman--Robbins.)} Write the mean gap as a $Q$-expectation against $W$. Since $\E_Q[W]=1$,
\[
\begin{aligned}
    a-b&=\E_P[T]-\E_Q[T]=\int T\,p-\int T\,q=\int T(p-q)\\
       &=\int q\,T\Big(\frac pq-1\Big)=\E_Q\big[T(W-1)\big].
\end{aligned}
\]
Because $\E_Q[W-1]=0$, this expectation is exactly a covariance: $a-b=\E_Q[T(W-1)]=\mathrm{Cov}_Q(T,W)$. By the Cauchy--Schwarz inequality for covariance,
\[
    (a-b)^2=\mathrm{Cov}_Q(T,W)^2\le\mathrm{Var}_Q(T)\,\mathrm{Var}_Q(W)=\mathrm{Var}_Q(T)\,\chi^2(P\|Q),
\]
and rearranging gives $\mathrm{Var}_Q(T)\ge (a-b)^2/\chi^2(P\|Q)$.

\emph{Cram\'er--Rao limit.} Put $Q=P_\theta$, $P=P_{\theta+h}$. Unbiasedness means $\E_{\theta'}[T]=\theta'$, so $a=\theta+h$, $b=\theta$, and $a-b=h$; thus $\mathrm{Var}_\theta(T)\ge h^2/\chi^2(P_{\theta+h}\|P_\theta)$. Taylor-expanding the density, $p_{\theta+h}=p_\theta+h\,\partial_\theta p_\theta+O(h^2)$, so $(p_{\theta+h}-p_\theta)^2=h^2(\partial_\theta p_\theta)^2+O(h^3)$ and
\[
    \chi^2(P_{\theta+h}\|P_\theta)=\int\frac{(p_{\theta+h}-p_\theta)^2}{p_\theta}
    = h^2\int\frac{(\partial_\theta p_\theta)^2}{p_\theta}+O(h^3)= h^2 I(\theta)+O(h^3),
\]
where $I(\theta)=\int (\partial_\theta\log p_\theta)^2\,p_\theta$ is the Fisher information. Letting $h\to0$,
\[
    \mathrm{Var}_\theta(T)\ge\frac{h^2}{h^2 I(\theta)+O(h^3)}\longrightarrow\frac{1}{I(\theta)}.
\]
So HCR is a finite-difference Cram\'er--Rao bound, with $\chi^2$ in the role of (squared) Fisher distance. $\qquad\blacksquare$
\item {\bf (CUBO.)} First, $p(x)$ is the $q$-average of the weight:
\[
    p(x)=\int p(x,z)\,dz=\int q(z\mid x)\,\frac{p(x,z)}{q(z\mid x)}\,dz=\E_q[w].
\]
Since the square is convex, Jensen gives $p(x)^2=(\E_q[w])^2\le\E_q[w^2]$; taking $\tfrac12\log$ (increasing) of both sides yields $\log p(x)\le\tfrac12\log\E_q[w^2]=\mathrm{CUBO}_2$.

\emph{Exact gap.} Factor the weight through the true posterior, $w=\frac{p(x,z)}{q(z\mid x)}=p(x)\,\frac{p(z\mid x)}{q(z\mid x)}$ (using $p(x,z)=p(x)\,p(z\mid x)$). Then
\[
\begin{aligned}
    \E_q[w^2]&=p(x)^2\,\E_q\!\Big[\Big(\frac{p(z\mid x)}{q(z\mid x)}\Big)^2\Big]
    =p(x)^2\int q(z\mid x)\,\frac{p(z\mid x)^2}{q(z\mid x)^2}\,dz\\
    &=p(x)^2\int\frac{p(z\mid x)^2}{q(z\mid x)}\,dz.
\end{aligned}
\]
By the foundation identity $\int\frac{p^2}{q}=1+\chi^2(p\|q)$ from part (a), the last integral is $1+\chi^2(p(z\mid x)\|q(z\mid x))$, so
\[
    \mathrm{CUBO}_2=\tfrac12\log\E_q[w^2]=\log p(x)+\tfrac12\log\!\big(1+\chi^2(p(z\mid x)\|q(z\mid x))\big).
\]
The gap $\tfrac12\log(1+\chi^2)\ge0$, and vanishes iff $\chi^2=0$, i.e.~$q(z\mid x)=p(z\mid x)$. Finally recall the ELBO is $\E_q[\log w]\le\log\E_q[w]=\log p(x)$ (Jensen, concavity of $\log$). Combining the two,
\[
    \mathrm{ELBO}\le\log p(x)\le\mathrm{CUBO}_2,
\]
bracketing the intractable log-evidence between two computable bounds. The ELBO gap is the \emph{reverse} KL $D(q(z\mid x)\|p(z\mid x))$ (mode-seeking), while the CUBO gap is driven by $\chi^2(p(z\mid x)\|q(z\mid x))$ (mass-covering). $\qquad\blacksquare$
\end{enumerate}
\end{solution}
\end{enumerate}

\begin{problem}
    \newpage~\newpage~\newpage~\newpage
\end{problem}
\end{document}
